\section*{Chapter \ref{ch:b1:aggregate-supply} --- Aggregate Supply and the Phillips Curve}
% Numbers recomputed in code/bachelor-1/25-aggregate-supply/exercises/solutions.py (tested).

\begin{solution}{exo:b1:aggregate-supply:sras-ex}
$Y = 10{,}000 + 2000 \times 0.05 = 10{,}100$.
\end{solution}

\begin{solution}{exo:b1:aggregate-supply:read}
1974 lies to the left of the crossing at $6.07\%$, where the line predicts rising
inflation ($2.853 - 0.47 \times 5.6 = 0.221$); the actual rise, $4.82$, is far above the line
--- the oil shock. 1982 lies to the right, where the line predicts a fall of
$0.47 \times 9.7 - 2.853 = 1.706$; the actual fall, $4.16$, is larger.
\end{solution}

\begin{solution}{exo:b1:aggregate-supply:predict}
$\pi = 3 - 0.5(4 - 5) = 3.5\%$.
\end{solution}

\begin{solution}{exo:b1:aggregate-supply:okun-ex}
Output is $2\%$ below its natural level.
\end{solution}

\begin{solution}{exo:b1:aggregate-supply:five}
Inflation rises by $0.5$ a year: $3.5$, $4$, $4.5$, $5$ and $5.5\%$.
\end{solution}

\begin{solution}{exo:b1:aggregate-supply:sacrifice}
$(10 - 4)/0.5 = 12$ point-years of excess unemployment; $2 \times 12 = 24$ points of a
year's output lost; a sacrifice ratio of $24/6 = 4$.
\end{solution}

\begin{solution}{exo:b1:aggregate-supply:estimate}
$\Delta\pi = 0$ at $u = 2.853/0.47 = 6.07\%$; $\beta = 0.47$.
\end{solution}

\begin{solution}{exo:b1:aggregate-supply:supply}
Holding $u = u^n$, inflation rises by 2 points and, with \omterm{def:b1:aggregate-supply:adaptive}{adaptive expectations},
stays 2 points higher. Holding inflation at its previous level requires
$u - u^n = v/\beta = 4$ points for one year: output falls by 4 times Okun's coefficient.
Policy chooses between the two, or a mix.
\end{solution}

\begin{solution}{exo:b1:aggregate-supply:demand}
In the short run AD shifts right: output rises above $\bar Y$ and prices rise by
less than $8\%$. As expected prices catch up, SRAS shifts up until output is
back at $\bar Y$; in the long run prices are $8\%$ higher and output unchanged.
\end{solution}

\begin{solution}{exo:b1:aggregate-supply:derive}
From \cref{prop:b1:aggregate-supply:sticky}, $P - P^e = (Y - \bar Y)/a$. With $P$ in logarithms
and last year's price level known, $P - P^e = \pi - \pi^e$. \omterm{def:b1:aggregate-supply:okun}{Okun's law} gives
$Y - \bar Y = -\omega\bar Y(u - u^n)$, so $\pi = \pi^e - (\omega\bar Y/a)(u - u^n)$: $\beta = \omega\bar Y/a$. A flatter SRAS
(larger $a$) gives a flatter Phillips curve.
\end{solution}

\begin{solution}{exo:b1:aggregate-supply:lambda}
With constant inflation $\pi$, $\pi = \lambda\pi + (1 - \lambda)\bar\pi - \beta(u - u^n)$, so
$u = u^n - (1 - \lambda)(\pi - \bar\pi)/\beta$: higher steady inflation buys lower unemployment
for ever if $\lambda < 1$, and not at all if $\lambda = 1$. In the 1970s expectations caught up
with inflation and the trade-off vanished: $\lambda$ was close to 1.
\end{solution}

\begin{solution}{exo:b1:aggregate-supply:refute}
\omterm{def:b1:aggregate-supply:adaptive}{Adaptive expectations} with $\lambda = 1$ predict the same relation between $\Delta\pi$ and $u$
in every period. A much flatter, weaker relation --- unemployment away from
$u^n$ without inflation accelerating --- refutes it. The slope of $-0.134$ with a
correlation of $-0.176$ over 1996--2018 points that way: expectations
anchored on a target ($\lambda < 1$). With 23 points it is suggestive, not decisive.
\end{solution}

\begin{solution}{pb:b1:aggregate-supply:kestria}
\textbf{1.} $5\%$.

\textbf{2.} $5\%$: inflation is constant only at the natural rate.

\textbf{3.} It rises by $0.5 \times 2 = 1$ point a year: $11$, $12$, $13$, $14$, $15\%$.

\textbf{4.} $\omega/\beta = 2/0.5 = 4$ points of output per point of inflation.

\textbf{5.} Any $u \ne 5$ changes inflation every year; only $u = 5$ keeps it constant,
at whatever level it is.

\textbf{6.} $4 = 10 - 0.5(u - 5)$: $u = 17\%$.

\textbf{7.} $5\%$: expectations are now 4.

\textbf{8.} $12$.

\textbf{9.} $2 \times 12 = 24$ points of a year's output.

\textbf{10.} Unemployment of $17\%$ for a year.

\textbf{11.} $9 = 10 - 0.5(u - 5)$ gives $u = 7\%$, and the same each year for six years.

\textbf{12.} $6 \times 2 = 12$.

\textbf{13.} $u = 11\%$ in each of the two years: again $12$ point-years.

\textbf{14.} Each year $u_t - 5 = (\pi_{t-1} - \pi_t)/0.5$; the sum telescopes to $(10 - 4)/0.5 = 12$, whatever
the path.

\textbf{15.} It contradicts it: across 65 disinflations, fast ones cost less per
point of inflation.

\textbf{16.} Expectations may respond to policy itself, not only to past
inflation; wage contracts may be renegotiated when a change of regime is
clear.

\textbf{17.} $10 - 4 = 6\%$.

\textbf{18.} $4 = 6 - 0.5(u - 5)$: $u = 9\%$, and $5\%$ afterwards.

\textbf{19.} $4$.

\textbf{20.} $12 - 4 = 8$ point-years.

\textbf{21.} $8$ points of output instead of $24$.

\textbf{22.} A record of meeting past targets, a central bank whose mandate and
independence make the target its own, and early action.

\textbf{23.} Inflation stays high, expected inflation goes back up, and the
next announcement is believed less.

\textbf{24.} Inflation fell from $9.1\%$ to $2.9\%$ with little rise in unemployment, as if
expectations had stayed anchored; the unwinding of supply shocks did part
of the work, so the episode is suggestive, not a test.

\textbf{25.} Under \omterm{def:b1:aggregate-supply:adaptive}{adaptive expectations} the cold-turkey and gradual paths cost
the same 12 point-years of unemployment; a credible announcement that
lowers expected inflation by 4 points saves 8.
\end{solution}
